\section{Carga}

Tipos de teste: carga, estresse, pico, resistência (\emph{soak}) e
capacidade.

Quatro sinais para observar:
\begin{enumerate}
\item \termo{vazão}{\emph{Throughput}: quantas requisições o sistema
  atende por unidade de tempo.};
\item latência em percentis (p95, p99; nunca só a média);
\item taxa de erro;
\item saturação.
\end{enumerate}

Prefira o modelo \emph{aberto} (taxa de chegadas) ao fechado, para não
cair na \termo{omissão coordenada}{\emph{Coordinated omission}: o gerador
de carga espera a resposta lenta antes de enviar a próxima requisição e,
assim, deixa de registrar a espera que usuários reais sofreriam. Origem da
expressão: Gil Tene.}. No Gatling, o modelo aberto usa \texttt{injectOpen}
(por exemplo \texttt{constant\-Users\-PerSec}) e o fechado, \texttt{injectClosed}
(\citetitle{doc:gatling-injection}).

Defina o critério antes de rodar. No Gatling, ele vira assertions que
quebram o build, como um limite de percentil de tempo de resposta e de taxa
de requisições com falha (\citetitle{doc:gatling-assertions};
\citetitle{artigo:threshold-test}).

Decida também o que fazer acima do limite: descartar com prioridade
(\termo{load shedding}{Descartar requisições de menor prioridade quando o
sistema está sobrecarregado, para proteger as importantes.}) é melhor do
que cair inteiro; veja os posts do Netflix TechBlog na tabela de artigos.
